\section{Sparse obstacles and local diffusion}\label{sec:sparse}

The return-energy problem requires two spatial estimates in each exposed
environment.  Mass initially on $T$ must spread through nearby holes,
and its pairing, after propagation, with the remainder $g=f-\1_T$ of
the hitting function must be bounded by the energy of that propagated
remainder.
Short connections near each starting point supply the first estimate;
flows through the rest of the lattice will supply the second.
We now derive the geometric conditions for both from the sparse pin-test
induction.  This section bounds the probability that these conditions
fail and proves the local Nash inequality.
Section~\ref{sec:flows} constructs the flows, and
Section~\ref{sec:leakage} uses both estimates to close the induction.

\subsection{The induction input and blocking between cuts}

Fix throughout
\begin{equation}\label{eq:sparse-exponents}
 \alpha=\frac1{100},\qquad \nu=\frac1{100d},\qquad a=\frac{11}{10}.
\end{equation}
All spatial distances in this and the following two sections are in the
supremum norm, and
\(B_h(x)=\{z\in\Z^d:|z-x|_\infty\le h\}\).
A finite collection of slot points with spatial coordinates \((x_i)_{i\in A}\)
is called \emph{sparse} if
\begin{equation}\label{eq:sparse-growth}
 \#\{i\in A:x_i\in B_h(z)\}\le L(1+h)^\alpha
 \qquad(z\in\Z^d,\ h\ge0).
\end{equation}
Repeated spatial coordinates are counted with their slot multiplicities.
The constant \(L\ge\max\{2d,\ell,4\}\) is fixed independently of all scales.
For a set of distinct sites, sparseness means the same inequality with one
point per site.  Every subset of a sparse collection is sparse.

Let \(0<p<1\), put
\begin{equation}\label{eq:sparse-scales}
 s=p^{-a},\qquad R=\lceil C_*s\rceil,
 \qquad \beta\ge5s,
\end{equation}
and choose a finite circular time grid \(\mathcal G\subset\R/\beta\Z\)
containing \(0\), whose successive gaps are at most \(p^{4d}\).
Here \(C_*\ge1\) is a fixed constant, to be chosen for the localization
estimate.  Write \(\mathcal V_{\mathcal G}=V_\Lambda\times\mathcal G\).
For \(D=D_0\cup F\), where \(F\subseteq\mathcal V_{\mathcal G}\), let
\(\mathsf I_p(D)\) denote the assertion
\begin{equation}\label{eq:induction-assertion}
 \mathbb P_D(A\subseteq\mathcal B_D)\le p^{|A|}
 \quad\text{for every nonempty sparse set }A\subseteq
 \mathcal V_{\mathcal G}\setminus D
 \text{ at a single grid time}.
\end{equation}
There are only finitely many grid points outside \(D_0\).
Thus one may prove these assertions by downward induction on the number
of pins, starting from the vacuous assertion when every grid point is
pinned.  In a step with pin set \(E\) and test \(T\), the set
\(P=E\cup T\) has strictly more pins.  In the present section we assume
\(\mathsf I_p(P)\); we do not assume \(\mathsf I_p(E)\).
Under this hypothesis we will seek a bound
\(r(E,T)\le\varepsilon(p)\) with \(\varepsilon(p)=o(p)\).
Proposition~\ref{prop:pin-test} then yields the required
\(p^{|T|}\) bound for the test under \(\mathbb P_E\), once
\(2\varepsilon(p)\le p\).

Expose \(\mathcal B_P\) under \(\mathbb P_P\) as in
Proposition~\ref{prop:exposure}.  A site is \emph{blocked} at time \(t\)
if none of its \(\ell\) slots is a hole.  Denote the blocked set by
\(\mathcal O_t\).  Extend the environment to \(\Z^d\) by making all
slots outside \(\Lambda\) available.  Boundary sites also have all slots
available because \(D_0\subseteq P\).  In particular \(\mathcal O_t\)
is finite and is disjoint from the inner boundary.

Sparse growth serves two purposes.  A blocked path crossing a large
spatial scale will supply a separated collection satisfying
\eqref{eq:sparse-growth}, so the restricted induction hypothesis can
still detect large obstacles.  Later, the same growth bound on the
test points will make overlaps between their decaying flows summable.
The component estimate below and the energy estimate in
Proposition~\ref{prop:flow-energy} implement these two uses.

At a grid time \(q\), the induction input gives
\begin{equation}\label{eq:grid-blocking}
 \mathbb P_P(A\subseteq\mathcal O_q)\le p^{|A|}
 \qquad\text{for every prescribed sparse set of distinct sites }A.
\end{equation}
Indeed, choose one fixed slot at each site.  If any chosen slot is pinned,
or any site is outside the volume, blocking is impossible.  Otherwise the
chosen slot points form an admissible test in \(\mathsf I_p(P)\).

\begin{lemma}[Blocking at a deterministic time]\label{lem:nongrid}
Assume \(\mathsf I_p(P)\), and let \(t\) be any deterministic time.
There is a constant \(C_0=C_0(d,\ell)\), uniform in \(p\le1\), the
volume, the period, the grid, and \(P\), such that, for every prescribed
finite set \(A\subset\Z^d\) of distinct sites,
\begin{align}
 \mathbb P_P(A\subseteq\mathcal O_t)&\le C_0p^{2d}
 &&\text{if }|A|=2d,                                      \label{eq:small-blocking}\\
 \mathbb P_P(A\subseteq\mathcal O_t)&\le(C_0\sqrt p)^{|A|}
 &&\text{if }A\text{ is a nonempty sparse set of distinct sites}.
                                                               \label{eq:sparse-blocking}
\end{align}
\end{lemma}

\begin{proof}
Let \(q\) be the preceding grid time and \(I=(q,t]\), interpreted
cyclically, so that \(|I|\le p^{4d}\).  On-site permutations do not
change the number of holes at a site.  Consequently a site blocked at
\(t\) was either blocked at \(q\) or has an incident inter-site mark
in \(I\).  This remains true when \(I\) crosses the time seam.
By Lemma~\ref{lem:factorial}, the expected number of marks in the
edge star of a prescribed site during \(I\) is at most
\(2d\ell^2|I|\).
If all \(2d\) prescribed sites are blocked at \(t\), either they were
all blocked at \(q\), or one of these finitely many edge stars contains
a mark.  Their set is sparse since \(L\ge2d\).  The union bound and
\eqref{eq:grid-blocking} give
\(p^{2d}+C p^{4d}\le Cp^{2d}\).

Now let \(n=|A|\) and assume \(A\) is sparse.  If every site of \(A\)
is blocked at \(t\), either at least \(\lceil n/2\rceil\) of its sites
were blocked at \(q\), or at least \(\lceil n/2\rceil\) have an
incident mark in \(I\).  For the first alternative, summing
\eqref{eq:grid-blocking} over subsets costs at most
\(2^n p^{\lceil n/2\rceil}\).
For the second, one of the two lattice parities supplies at least
\(k=\lceil n/4\rceil\) sites with incident marks.  Edge stars of
distinct sites of the same parity are disjoint.  If \(N_z\) counts
marks in the star of \(z\) during \(I\), then for any such prescribed
\(k\)-set \(A'\), the product \(\prod_{z\in A'}N_z\) counts
choices of \(k\) distinct marks.  The factorial bound therefore yields
\[
 \mathbb P_P(N_z\ge1\text{ for all }z\in A')
 \le \mathbb E_P\prod_{z\in A'}N_z
 \le (2d\ell^2|I|)^k.
\]
Summing over at most \(2^{n+1}\) possible choices bounds the second
alternative by \(2^{n+1}(C p^{4d})^{\lceil n/4\rceil}\).
As \(4d\lceil n/4\rceil\ge dn\ge n/2\), both alternatives are
bounded together by \((C_0\sqrt p)^n\), after enlarging \(C_0\).
These are separate event bounds combined by a union bound; no
independence between old blocking and intervening marks is required.
\end{proof}

\subsection{Admissible binary witnesses and component tails}

Join two blocked sites when their supremum distance is at most \(8\).
All components and paths of blocked sites below use this adjacency.
For \(z\in\mathcal O_t\), let \(\mathcal C_t(z)\) be its component
and set
\[
 \operatorname{rad}_z\mathcal C_t(z)
 =\max_{w\in\mathcal C_t(z)}|w-z|_\infty.
\]
The next lemma gives a tail uniform even as the bootstrap parameter is
subsequently decreased.

\begin{lemma}[Uniform component tail]\label{lem:component-tail}
There are \(p_0>0\) and \(c,C,\gamma>0\), depending only on
\(d,\ell\) and the fixed sparse-growth choices, such that whenever
\(0<p\le p_0\) and \(\mathsf I_p(P)\) holds,
\begin{equation}\label{eq:component-tail}
 \mathbb P_P\bigl(z\in\mathcal O_t,\
      \operatorname{rad}_z\mathcal C_t(z)\ge h\bigr)
 \le C\exp(-ch^\gamma)
 \qquad(z\in\Z^d,\ h\ge1)
\end{equation}
at every deterministic time \(t\).  The constants \(c,C,\gamma\)
are fixed once \(p_0\) is fixed and work simultaneously for all
\(p\le p_0\).
\end{lemma}

\begin{proof}
Fix an integer \(b>100\) so large that \(b^\alpha>4\), and put
\(h_j=16b^j\).  We first describe exactly which site lists will be
used as witnesses.  Their separated binary-tree structure is also used
in multiscale cascading-event arguments; see, for example,
\cite[Section~3]{Sznitman2012}.  Here the leaf probabilities will come
directly from \eqref{eq:sparse-blocking}.  A depth-\(j\) witness rooted at \(z\) is an
assignment of sites \(x_w\) to the vertices \(w\) of a rooted ordered
binary tree of depth \(j\), with \(x_\varnothing=z\).  If a vertex
\(w\) has \(k\ge1\) levels below it, require
\begin{equation}\label{eq:witness-admissibility}
 x_{w0}=x_w,\qquad x_{w1}\in B_{h_k+8}(x_w),\qquad
 \operatorname{dist}_\infty(\mathcal L_{w0},\mathcal L_{w1})
       \ge h_k/4,
\end{equation}
where \(\mathcal L_{wa}\) is the set of leaf sites descended from
child \(wa\).  Call such an assignment \emph{admissible}.
Its leaves are distinct by the separation condition.  Depth zero has
just the prescribed root as its leaf.

A blocked crossing from \(z\) to distance at least \(h_j\) produces
an admissible witness all of whose leaves are blocked.  For \(j=0\)
use the root alone.  For \(j\ge1\), truncate the crossing at its
first exit to distance \(h_j\); it is
contained in \(B_{h_j+8}(z)\).  One child crossing starts at \(z\).
The other starts at the first site whose distance from \(z\) is at
least \(h_j/2\), hence at distance at most \(h_j/2+8\).
The remaining part of the parent crossing reaches distance at least
\(h_j/2-8\) from this second start.  Each child can thus be truncated
at its first exit to distance \(h_{j-1}\), and is contained in its
start-centered ball of radius \(h_{j-1}+8\).  The distance between
these two balls is at least
\[
 h_j/2-2(h_{j-1}+8)\ge h_j/4.
\]
Recurse on the two truncated child crossings.  Every descendant leaf
remains in its child's truncated crossing, so all the separation
conditions in \eqref{eq:witness-admissibility} hold.

Every admissible witness has a sparse leaf set.  Indeed, for
\(1\le k\le j\), a ball of diameter strictly less than \(h_k/4\)
meets at most one child at every split of scale at least \(h_k\).
It therefore contains at most \(2^{k-1}\) leaves.  If
\(h_{k-1}/8\le h<h_k/8\), this bound is at most
\(L(1+h)^\alpha\), since
\[
 L(1+h)^\alpha\ge L(2b^{k-1})^\alpha
       \ge 2^{k-1}.
\]
For \(h<h_0/8\) there is at most one leaf, and for
\(h\ge h_j/8\) the same calculation allows all \(2^j\) leaves.
The depth-zero case is immediate.  This verifies
\eqref{eq:sparse-growth} for all centers and real radii, including
the scale endpoints.

Let \(A_j\) be an upper bound on the number of admissible assignments
rooted at any prescribed site.  At the top split the first child root
is fixed and the second has at most \((C_dh_j)^d\) choices.
Ignoring restrictions between the two child assignments gives the
numerical overcount
\[
 A_0=1,\qquad A_j\le(C_dh_j)^d A_{j-1}^2.
\]
In particular
\[
 2^{-j}\log A_j
 \le d\sum_{k=1}^j2^{-k}\log(C_dh_k)
 \le \log C_b<\infty,
 \qquad A_j\le C_b^{2^j}.
\]
Only admissible assignments enter the probability union bound.  The
unrestricted product used to bound their number need not itself
consist of sparse witnesses.

For each admissible assignment, Lemma~\ref{lem:nongrid} bounds the
probability that all its leaves are blocked by
\((C_0\sqrt p)^{2^j}\).  Choose \(p_0\le1\) once and for all so
that \(C_bC_0\sqrt{p_0}\le e^{-2}\).  For every \(p\le p_0\),
\begin{equation}\label{eq:discrete-component-tail}
 \mathbb P_P\bigl(z\in\mathcal O_t,\
      \operatorname{rad}_z\mathcal C_t(z)\ge h_j\bigr)
 \le \exp(-2\cdot2^j).
\end{equation}
For \(h\ge16\), choose \(j\ge0\) with
\(h_j\le h<h_{j+1}\).  Then
\(2^j\ge(h/(16b))^\gamma\), where \(\gamma=\log_b2>0\).
This proves \eqref{eq:component-tail}; increasing \(C\) covers
\(1\le h<16\).  No constant in this last step changes when \(p\)
is decreased further.
\end{proof}

\subsection{Regular anchors}

We now separate the geometry near the starting point of a localized
mass from the geometry farther away.  Nearby, fewer than $2d$ blocked
sites in each fixed-radius ball will preserve short lattice connections;
these connections transfer the lattice Nash inequality to the holes.
Farther away, a sublinear bound on component radii will control the
displacement of the routes used for flows to infinity in
Section~\ref{sec:flows}.  The first condition gives local diffusion,
and the second permits the global energy comparison.
Let \(M_*\) be a fixed positive integer.  Its required size will be
specified below, independently of \(p\).

\begin{definition}[Regular anchor]\label{def:regular-anchor}
For a spatial anchor \(x\in\Z^d\), call time \(t\)
\emph{regular for \(x\)} at scale \(R\) if both conditions hold:
\begin{enumerate}
 \item For every \(z\in B_{11R}(x)\),
       \(\#(\mathcal O_t\cap B_{M_*}(z))<2d\).
 \item For every \(z\in\mathcal O_t\setminus B_{10R}(x)\),
       \(\operatorname{rad}_z\mathcal C_t(z)
            \le(1+|z-x|_\infty)^\nu\).
\end{enumerate}
For \(i\in T\), ``regular for \(i\)'' means regular for its spatial
coordinate \(x_i\).
\end{definition}

\begin{proposition}[Probability of irregularity]\label{prop:regularity}
Fix \(M_*\).  Under \(\mathsf I_p(P)\), for all \(0<p\le p_0\),
all deterministic times \(t\), and all spatial anchors \(x\),
\begin{equation}\label{eq:regularity-failure}
 \mathbb P_P(t\text{ is not regular for }x)
 \le \delta:=C R^d p^{2d}+C\exp(-cR^\xi),
 \qquad \xi=\nu\gamma>0.
\end{equation}
Here \(R\ge1\), and the constants may depend on \(M_*\) but are
uniform in \(p\le p_0\), volume, period, grid, pin set and anchor.
\end{proposition}

\begin{proof}
There are at most \(C R^d\) centers in \(B_{11R}(x)\), and a ball
of radius \(M_*\) has a fixed number of \(2d\)-element subsets.
The union bound and \eqref{eq:small-blocking} therefore bound failure
of the first condition by \(C R^d p^{2d}\).
By Lemma~\ref{lem:component-tail}, failure of the second has probability
at most
\[
 C\sum_{z:|z-x|_\infty>10R}
   \exp\bigl(-c(1+|z-x|_\infty)^{\nu\gamma}\bigr)
 \le C\exp(-c'R^{\nu\gamma}).
\]
For the last inequality, count at most \(C(1+n)^{d-1}\) sites at
distance \(n\) and absorb this polynomial and the tail summation
into half of the exponential.  Since \(\nu\gamma>0\) is fixed,
the resulting constants are uniform.  Add the two bounds.
\end{proof}

We will use \eqref{eq:regularity-failure} only at prescribed times and
after integration in time.  It asserts neither independence between
slices nor regularity simultaneously at every time.

\subsection{Local connections and the Nash inequality}

The remaining statements in this section are deterministic.  At a
fixed time let \(\mathcal H_t\) be the graph whose vertices are the
available slots on the extended lattice and whose edges join every
pair of available slots above nearest-neighbor sites.  For a real
function \(v\) on its vertices define
\begin{equation}\label{eq:slot-energy}
 \mathcal D_t(v)=\frac12
 \sum_{\{i,j\}\in E(\mathcal H_t)}(v_i-v_j)^2.
\end{equation}
Edges in this sum are unoriented and counted once.  Norms use counting
measure on the available slots.  When a function describes killed
mass, its values at the physical boundary and outside the volume are
zero.  Thus \eqref{eq:slot-energy} includes the edges into the
zero-valued boundary, as required by Dirichlet killing.

\begin{lemma}[Uniform local connectivity]\label{lem:local-connectivity}
Let \(d\ge2\) and \(K\ge1\) be fixed.  There is an integer
\(\rho=\rho(d,K)\ge K+1\) such that the following holds.
If \(u,v\in\Z^d\), \(|u-v|_\infty\le K\), and
\(F\subset\Z^d\) satisfies
\[
 u,v\notin F,\qquad \#(F\cap B_\rho(u))<2d,
\]
then a nearest-neighbor path from \(u\) to \(v\) avoids \(F\),
lies in \(B_\rho(u)\), and has length at most \((2\rho+1)^d\).
\end{lemma}

\begin{proof}
First suppose \(F\) is a set of fewer than \(2d\) sites in the
whole lattice.  The complement of a box containing \(F\) is connected
when \(d\ge2\), so \(\Z^d\setminus F\) has only one possible
infinite component.  It has no finite component: from a vertex of
such a component the first exit along each of the \(2d\) positive
and negative coordinate rays must lie in \(F\), giving \(2d\)
distinct deleted sites.  Hence \(\Z^d\setminus F\) is connected.

The confinement radius can be uniform when the two endpoints have
bounded separation.  Otherwise, after translating the first endpoint
to \(0\), there would be sets \(F_n\) with \(|F_n|<2d\) and
surviving endpoints \(v_n\in B_K(0)\), admitting no path inside
\(B_n(0)\).  Pass to a subsequence with \(v_n=v\) and with the
indicators of \(F_n\) converging pointwise on \(\Z^d\).  The limiting
set \(F_\infty\) has fewer than \(2d\) sites and avoids both
endpoints.  By the preceding paragraph there is a finite path from
\(0\) to \(v\) avoiding \(F_\infty\).  For all sufficiently large
\(n\), this path lies in \(B_n(0)\) and avoids \(F_n\), a
contradiction.  Choose a uniform radius \(\rho\ge K+1\).

For the asserted local version apply this conclusion to
\(F'=F\cap B_\rho(u)\), a set of fewer than \(2d\) sites in the
whole lattice.  The resulting path lies in \(B_\rho(u)\), where
\(F'\) and \(F\) agree.  Erasing loops gives the claimed length bound.
\end{proof}

We may now make the local radius choice without reference to \(p\).
Fix
\begin{equation}\label{eq:local-radius-choice}
 K_{\mathrm{conn}}=8(2d-2)+12,\qquad
 M_*\ge\max\{8(2d),\rho(d,K_{\mathrm{conn}})+2\}.
\end{equation}
We henceforth require \(R\) to exceed these fixed local radii, which
is achieved by decreasing the final \(p\).  At a regular anchor,
Lemma~\ref{lem:local-connectivity} is now available for surviving
sites within separation \(K_{\mathrm{conn}}\), whenever its first
endpoint lies in \(B_{11R}(x)\).  In particular, it applies to every
short connection used below.  This choice also serves the local
source connections in Section~\ref{sec:flows}.

The next estimate uses Nash's norm--energy inequality and its diffusion
mechanism~\cite[Part~I]{Nash1958}.  The transfer to the available slots
uses path comparison of Dirichlet forms, a method developed
systematically by Diaconis and Saloff-Coste
\cite[Theorem~2.1]{DiaconisSaloffCoste1993}.  We supply the comparison
for the present obstacles and zero extension explicitly.

\begin{proposition}[Local Nash inequality]\label{prop:local-nash}
Fix \(d\ge3\), \(\ell\), and \(M_*\) as in
\eqref{eq:local-radius-choice}.  Let \(\mathcal H_t\) be an available-slot
graph with at least one and at most \(\ell\) slots at every unblocked
site, and suppose \(t\) is regular for \(x\) at scale
\(R\ge M_*+K_{\mathrm{conn}}+3\).
Every nonnegative function \(v\) supported on slots above \(B_R(x)\)
satisfies
\begin{equation}\label{eq:local-nash}
 \|v\|_2^{2+4/d}\le C\mathcal D_t(v)\|v\|_1^{4/d},
\end{equation}
where \(C\) depends only on \(d,\ell\) and the fixed local radii.
\end{proposition}

\begin{proof}
We transfer \(v\) to a lattice function while controlling energy, then
prove the lattice inequality by averaging.  At every surviving site
\(z\), choose an available slot maximizing \(v\), and let \(V(z)\)
be this maximum.  Every blocked site in \(B_{R+3}(x)\) has a surviving
nearest neighbor: otherwise its \(2d\) neighbors alone would violate
the first regularity condition in the ball centered there.  Choose
one such neighbor \(\phi(z)\) and define
\[
 w(z)=
 \begin{cases}
   V(z),&z\notin\mathcal O_t,\\
   V(\phi(z)),&z\in\mathcal O_t\cap B_{R+3}(x),\\
   0,&z\in\mathcal O_t\setminus B_{R+3}(x).
 \end{cases}
\]
The function \(w\) is supported in \(B_{R+1}(x)\).  Since a site's
value is copied only by its nearest neighbors,
\begin{equation}\label{eq:nash-norm-transfer}
 \|v\|_2^2\le\ell\|w\|_2^2,
 \qquad \|w\|_1\le(2d+1)\|v\|_1.
\end{equation}

For a lattice edge \(\{z,z'\}\) with a nonzero endpoint value,
both endpoints belong to \(B_{R+2}(x)\).  Each of \(w(z),w(z')\)
is the value at a chosen maximizing slot over a surviving site,
say \(u,u'\), at distance at most one from the corresponding
endpoint.  Thus \(|u-u'|_\infty\le3\) and \(u,u'\in B_{R+3}(x)\).
Lemma~\ref{lem:local-connectivity} joins them through surviving sites
within a fixed radius of \(u\), with a fixed length bound.  Lift
this path using the chosen maximizing slot at every site.  Every
lifted edge is available because all pairs of available slots across
a lattice edge are joined.  If \(u=u'\), the difference is zero.
Otherwise telescoping and Cauchy--Schwarz show that
\((w(z)-w(z'))^2\) is at most a fixed constant times the sum of
\((v_i-v_j)^2\) along the lifted path.
For any fixed slot edge, only a bounded number of the original
lattice edges can have a path using it: all of their endpoints lie
within a fixed radius of that edge.  Hence, summing,
\begin{equation}\label{eq:nash-energy-transfer}
 \sum_{\{z,z'\}:|z-z'|_1=1}(w(z)-w(z'))^2
 \le C\mathcal D_t(v).
\end{equation}

For completeness, write \(\mathcal E_{\Z^d}(w)\) for the sum on
the left of \eqref{eq:nash-energy-transfer}.  For an integer
\(k\ge1\), let \(A_kw\) be the average of translates of \(w\)
by the points of \(\{0,\ldots,k-1\}^d\).  Telescoping a translate
along coordinate directions, followed by Cauchy--Schwarz, gives
\[
 \|w-A_kw\|_2^2\le C_d k^2\mathcal E_{\Z^d}(w).
\]
The averaging kernel has counting \(2\)-norm \(k^{-d/2}\);
the convolution inequality therefore gives
\(\|A_kw\|_2^2\le k^{-d}\|w\|_1^2\).  Consequently
\[
 \|w\|_2^2\le C_d k^2\mathcal E_{\Z^d}(w)
                 +2k^{-d}\|w\|_1^2.
\]
If \(w\ne0\), choose
\(k=\lceil(4\|w\|_1^2/\|w\|_2^2)^{1/d}\rceil\).
The second term is at most \(\|w\|_2^2/2\), and
\(\|w\|_1^2/\|w\|_2^2\ge1\) bounds the effect of rounding.
Absorption yields
\[
 \|w\|_2^{2+4/d}
 \le C_d\mathcal E_{\Z^d}(w)\|w\|_1^{4/d}.
\]
For \(w=0\) the inequality is immediate.  Finally use
\eqref{eq:nash-norm-transfer} and \eqref{eq:nash-energy-transfer}
to obtain \eqref{eq:local-nash}.
\end{proof}

The order of choices is now explicit.  First fix \(d,\ell\), the
exponents in \eqref{eq:sparse-exponents}, the witness scale ratio
\(b\), and \(L\).  Next fix the preliminary \(p_0\) and the uniform
tail constants in Lemma~\ref{lem:component-tail}.  Fix the connectivity
radii and \(M_*\) by \eqref{eq:local-radius-choice}; this fixes the
constants in Proposition~\ref{prop:regularity} and
Proposition~\ref{prop:local-nash}.  Choose \(C_*\) for localization,
and only then decrease \(p\le p_0\) so that \(R\) exceeds the fixed
radii and the leakage error is small enough to close the induction.
Thus no tail constant deteriorates as the final \(p\) is decreased.
The next section uses the same regularity conditions to construct
flows to infinity, complementing the local diffusion bound proved here.
